\documentclass[11pt]{article}
\usepackage{preamble}
\setlength{\parskip}{4pt}

\begin{document}

\daybanner{AP Statistics \textperiodcentered\ Unit 5 \textperiodcentered\ Topic 5.1 \textperiodcentered\ B: 2027-02-05 \textperiodcentered\ E: 2027-03-31 \textperiodcentered\ TEACHER KEY}{One Plot, Two Reporting Strategies}

\sectionbanner{CED TETHER}

{\small
\begin{itemize}[leftmargin=1.5em,itemsep=2pt,topsep=2pt]
  \item \textbf{Skill 2.B} --- Construct numerical or graphical representations of distributions.
  \item \textbf{Skill 2.A} --- Describe data presented numerically or graphically.
  \item \textbf{EU UNC-1} --- Graphical representations and statistics allow us to identify and represent key features of data.
  \item \textbf{LO UNC-1.S} --- Represent bivariate quantitative data using scatterplots. [Skill 2.B]
  \item \textbf{EK UNC-1.S.1} --- A bivariate quantitative data set consists of observations of two different quantitative variables made on individuals in a sample or population.
  \item \textbf{EK UNC-1.S.2} --- A scatterplot shows two numeric values for each observation, one corresponding to the value on the x-axis and one corresponding to the value on the y-axis.
  \item \textbf{EK UNC-1.S.3} --- An explanatory variable is a variable whose values are used to explain or predict corresponding values for the response variable.
  \item \textbf{EU DAT-1} --- Regression models may allow us to predict responses to changes in an explanatory variable.
  \item \textbf{LO DAT-1.A} --- Describe the characteristics of a scatter plot. [Skill 2.A]
  \item \textbf{EK DAT-1.A.1} --- A description of a scatter plot includes form, direction, strength, and unusual features.
  \item \textbf{EK DAT-1.A.2} --- The direction of the association shown in a scatterplot, if any, can be described as positive or negative.
  \item \textbf{EK DAT-1.A.3} --- A positive association means that as values of one variable increase, the values of the other variable tend to increase. A negative association means that as values of one variable increase, values of the other variable tend to decrease.
  \item \textbf{EK DAT-1.A.4} --- The form of the association shown in a scatterplot, if any, can be described as linear or non-linear to varying degrees.
  \item \textbf{EK DAT-1.A.5} --- The strength of the association is how closely the individual points follow a specific pattern, e.g., linear, and can be shown in a scatterplot. Strength can be described as strong, moderate, or weak.
  \item \textbf{EK DAT-1.A.6} --- Unusual features of a scatter plot include clusters of points or points with relatively large discrepancies between the value of the response variable and a predicted value for the response variable.
\end{itemize}
}
\textbf{Essential question:} How should several scatterplot features be weighed when descriptions conflict?\par
\textbf{DOK-3 skill:} 4.B \textperiodcentered\ \textbf{FRQ pattern:} {\small\texttt{linear-\allowbreak{}label-\allowbreak{}vs-\allowbreak{}curved-\allowbreak{}pattern-\allowbreak{}description-\allowbreak{}and-\allowbreak{}consequence}}\par
\textbf{DOK rationale:} Part (c) weighs competing claims using the day's numerical or graphical evidence and explains a limitation or consequence; the judgment requires linked reasoning beyond a calculation.\par

\frameworkphaseheader{Do Now $\rightarrow$ video follow-along $\rightarrow$ finish and turn in}{1 $\rightarrow$ 3}{45}{%
\begin{itemize}[leftmargin=1.5em,itemsep=2pt,topsep=2pt]
  \item Show the board at the bell and collect a one-sentence first take without confirming a claim.
  \item Run the named video follow-along, then allow about 10 minutes for parts (a)--(c).
  \item Ask for evidence linked to a claim. Collect the sheet and hand-score part (c) with E/P/I.
\end{itemize}
}{%
\begin{itemize}[leftmargin=1.5em,itemsep=2pt,topsep=2pt]
  \item Commit to a reporting strategy, then complete the follow-along.
  \item Finish (a)--(c), revise the first take, and turn in.
\end{itemize}
}{%
\begin{itemize}[leftmargin=1.5em,itemsep=2pt,topsep=2pt]
  \item Which number or visible feature supports your claim?
  \item Why does the other claim go beyond that evidence?
  \item What would you tell someone who chose the other claim?
\end{itemize}
}{Circulate and ask students to connect their choice to evidence. Score part (c) using the three required reasoning elements; do not require dropped extension work.}

\begin{callout}{\IconWarn\ WATCH FOR}{calloutred}
\begin{itemize}[leftmargin=1.5em,itemsep=2pt,topsep=2pt]
  \item A choice with copied numbers but no explanation of how they support it.
  \item Treating a model or average as a guaranteed individual outcome.
\end{itemize}
\end{callout}

\sectionbanner{SCORING PART (c) --- E / P / I}

\textbf{Expected elements}
\begin{itemize}[leftmargin=1.5em,itemsep=2pt,topsep=2pt]
  \item \textbf{reasoning-1} (required) --- Supports Ravi with direction, form, and strength in context
  \item \textbf{reasoning-2} (required) --- Uses 170 versus 15 mL and identifies the low 50-minute batch
  \item \textbf{reasoning-3} (required) --- Explains how a straight line could overstate later growth
\end{itemize}

{\small\renewcommand{\arraystretch}{1.25}
\noindent\begin{tabular}{|p{0.5in}|p{5.9in}|}
\hline
\textbf{E} & Includes all three required reasoning elements above, with correct contextual evidence. \\ \hline
\textbf{P} & Includes at least one substantive correct reasoning link above, but omits or misstates another required element. \\ \hline
\textbf{I} & Includes no substantive correct reasoning link above; an unsupported choice or copied number alone is insufficient. \\ \hline
\end{tabular}}

\textbf{Common mistakes}
\begin{itemize}[leftmargin=1.5em,itemsep=2pt,topsep=2pt]
  \item Calling every positive association linear
  \item Letting one point determine the entire form
  \item Omitting context, strength, or the practical consequence
\end{itemize}

\clearpage
\focusbanner{TODAY'S PROBLEM --- annotated}

Suppose a test kitchen measures dough volume for ten batches at assigned times after mixing. The scatterplot shows the paired measurements.\par\smallskip \textbf{Jules:} Use a straight-line summary because volume generally increases with time. Treat the 50-minute batch as ordinary variation.\par \textbf{Ravi:} Decide the form from how the rate changes, and report the 50-minute batch separately if it departs from nearby values.\par Selected measurements (minutes, mL): $(0,420)$, $(20,590)$, $(70,770)$, $(90,785)$.\par
\begin{center}
\begin{tikzpicture}
\begin{axis}[
  width=0.56\linewidth,
  height=1.51in,
  xmin=0, xmax=90, ymin=380, ymax=820,
  xlabel={Minutes after mixing}, ylabel={Dough volume (mL)}, grid=major,
  tick label style={font=\small}, label style={font=\small}
]
\addplot[only marks, mark=*, mark size=1.4pt, color=dayblue] coordinates {
  (0,420)
  (10,510)
  (20,590)
  (30,650)
  (40,695)
  (50,600)
  (60,750)
  (70,770)
  (80,780)
  (90,785)
};
\end{axis}
\end{tikzpicture}
\end{center}
\begin{firsttakebox}{FIRST TAKE (5 min)}
Which reporting strategy seems more defensible, Jules's or Ravi's? Cite one graph feature.\par
\answer{Not graded; students revise after coordinating all features.}
\end{firsttakebox}

\textit{Rules callout ``BUILD A COMPLETE SCATTERPLOT DESCRIPTION (keep on the page)'' is printed on the student sheet.}\par\medskip

\dokbadge{1}~\textbf{(a)} Name the explanatory and response variables, with units.\par
\answer{Time after mixing (minutes) is explanatory; dough volume (mL) is the response.}

\dokbadge{2}~\textbf{(b)} Find the volume increase from 0 to 20 minutes and from 70 to 90 minutes. Is the 50-minute batch above or below the nearby batches?\par
\answer{The increases are $590-420=170$ mL and $785-770=15$ mL. The 50-minute batch at 600 mL is below the nearby 40- and 60-minute batches (695 and 750 mL).}

\dokbadge{3}~\textbf{(c)} Whose reporting strategy is supported? Describe the direction, form, strength, and unusual point in context, using your two increases. Explain how a straight-line summary could mislead someone predicting later growth. Write 3--4 sentences.\par
\answer{Ravi is supported: the association is fairly strong and positive, but nonlinear because most volumes rise and then level off. The 170-mL early increase versus 15 mL late supports the curve, and the 50-minute batch is unusually low compared with nearby batches. A straight-line summary could overstate later growth by suggesting that the early rate continues.}

\begin{summaryexitbox}{EXIT}
One thing the video changed about my first take: \hrulefill
\end{summaryexitbox}

\end{document}
